\chapter{Direct Antiglobulin Test (Coombs Test)}

\section*{What Is This Test?}
The direct antiglobulin test looks for immunoglobulin or complement attached to a patient's red blood cells in the circulation. It is a key laboratory test for immune-mediated hemolysis.

\section*{Alternative Names}
\begin{itemize}
  \item DAT
  \item Direct Coombs Test
  \item Direct Antiglobulin Reaction
\end{itemize}
\section*{Why This Test Is Ordered}
It is used when autoimmune hemolytic anemia, hemolytic transfusion reactions, hemolytic disease of the fetus or newborn, or selected drug-related hemolysis is suspected. It helps separate immune red-cell destruction from nonimmune causes. Still, it must be interpreted along with evidence of hemolysis.

\section*{How This Test Is Performed}
The patient's red cells are washed and then exposed to antihuman globulin reagent. Agglutination indicates that IgG, complement, or both are present on the red-cell surface. Monospecific reagents can tell IgG coating apart from complement coating.

\section*{Preparation, Risks, and Limitations}
No fasting is required. Risks are limited to venipuncture. A positive result does not prove that hemolysis is happening. A negative result does not exclude every immune hemolytic process. Recent transfusion, medicines, and technical factors can affect interpretation.

\section*{Understanding Results}
A positive DAT together with anemia, high reticulocytes, high indirect bilirubin, high LDH, and low haptoglobin supports immune hemolysis. The clinical pattern and transfusion history decide whether the likely cause is autoimmune, transfusion-related, neonatal, or drug-associated.

\section*{References}
British Society for Haematology guidance on autoimmune hemolytic anemia and pretransfusion testing.

\clearpage
\section*{Understanding Results and Follow-up}
The laboratory report should identify the specimen, method, units, reference interval or decision limit, and whether the result is positive, negative, abnormal, or inconclusive. Reference intervals vary with age, sex, pregnancy, specimen type, assay, and laboratory; a result outside a range is not automatically a diagnosis, and a result within a range does not always exclude disease. Discuss unexpected results with the ordering clinician, who can decide whether repeat or confirmatory testing, treatment, or referral is appropriate. Do not change medicines or delay urgent care based on an isolated result.


\section*{Regulatory Status and Authorization Date}
This chapter describes a test category, not a specific commercial product. FDA, CDSCO, EU IVDR, and MHRA approval or registration dates therefore cannot be assigned without the assay manufacturer, product name, instrument, and intended use. For laboratory-developed tests, an individual agency approval date may not exist. See the Preface for the interpretation of regulatory status.

\clearpage
